\lecture{6}{Thu 05 Feb 2026 12:28}{Complex furries}

\begin{remark}
	I finally have good motivation for why complexifying the tangent space is tensoring with $\mathbb{C} $. Let $v\in T_pM$ and $z=x+iy\in\mathbb{C} $. Then $v\otimes z\in T_p^\mathbb{C} M$ is equivalent to, by bilinearity,
	\[
		v\otimes (x+iy)=v\otimes x + v\otimes iy
	.\]
	Of course by $\mathbb{R} $-linearity of $\otimes _\mathbb{R} $, this is equal to
	\[
		=x(v\otimes 1)+y(v\otimes i)
	.\]
	We can choose to write $v\otimes 1=v$ and $v\otimes i=iv$, and this really does behave like normal multiplication since normal multiplication is also bilinear. Thus all elements of $T_p^\mathbb{C} M$ can be written of the form $u+iv$ for some $u,v\in T_pM$; in our example above, $u=xv$ and $v=yv$ (abusing notation very badly). This is required to have $\partial _z=\frac{1}{2} (\partial _x -i \partial _y)$ in the tangent space.
\end{remark}

\begin{proposition}\label{???}
	Let $X$ be a compact complex connected submanifold, and $f : X \to \mathbb{C} $ a holomorphic function. Then $f$ is constant.
\end{proposition}
\begin{proof}
	Easy to show $f$ is constant on connected components by passing to coordinates.
\end{proof}
\begin{corollary}\label{????}
	If $X\subseteq \mathbb{C} ^n$ is a compact complex connected submanifold, then $X$ is a point.
\end{corollary}
\begin{proof}
	Take the holomorphic function
	\[\begin{tikzcd}[row sep = 2.5em, column sep = 2.5em]
	X \ar[hook, r]  & \mathbb{C} ^n \ar[" \pi _i ", r]  & \mathbb{C} .
	\end{tikzcd}\]
	This must be constant by \cref{???}.
\end{proof}

Since looking for compact complex connected subsets of $\mathbb{C} ^n$, we look to projective space for something like Whittney's theorem for embedding complex manifolds into $\mathbb{C} ^n$. This still is not perfect, but is worth considering.

\section{The projective space}

Recall that projective space $k\mathbb{P} ^n$ over a field $k$ is defined as the set of lines through the origin in $k ^{n+1} \setminus \left\{ 0 \right\} $, topologized in the obvious way:
\[
	k\mathbb{P} ^n \coloneqq \frac{k^{n+1}\setminus \left\{ 0 \right\} }{ \left\langle \lambda x \sim x \mid x\in k^{n+1} ,\lambda \in k\setminus \left\{ 0 \right\}  \right\rangle 
}.\]
Given $\pi $ the canonical projection and $L$ any line in $\mathbb{C} ^{n+1} $, we have $\pi ^{-1} (L)\cong L\setminus \left\{ 0 \right\} \cong \mathbb{C} ^*$. We can then show that $\mathbb{C} ^{n+1} \setminus 0$ is a locally trivial $\mathbb{C} ^*$-bundle over $\mathbb{C} \mathbb{P} ^n$. Recall that the atlas is formed by sets
\[
	U_i \coloneqq \left\{ [c_0 : \cdots : c_n]\in\mathbb{C} \mathbb{P} ^n \mid c_i \neq 0 \right\} 
.\]
This is well-defined. The coordinates simply drop $c_i$, since we can always take it to be normalized at 1.

\begin{remark}
	We work in the quotient topology on $\mathbb{C} \mathbb{P} ^n$, which is not the same as the usual topologies we would use in algebraic geometry. However many of the results translate between the two, due to some work by Serre.
\end{remark}

We equip $\mathbb{C} ^{n+1} $ with the Hermitian inner product $\left\langle z,w \right\rangle =\sum_{i} z_i \overline{w} _i$, the \emph{opposite} of the physics convention.

The fibration
\[\begin{tikzcd}[row sep = 2.5em, column sep = 2.5em]
S^1 \ar[hook, r]  & S^{n+1} \ar[r]  & \mathbb{C} \mathbb{P} ^n
\end{tikzcd}\]
is called the Hopf fibration.

Complex projective space is a CW complex. We can define $\mathbb{C} \mathbb{P} ^1 \coloneqq \mathbb{C} ^1 \amalg \mathbb{C} ^0$ with $\mathbb{C} ^0$ the point at $\infty $. If we take a chart $U_0 \subseteq \mathbb{C} \mathbb{P} ^2$, then $\mathbb{C} \mathbb{P} ^2\setminus U_0$ is just the set $\left\{ [0 : z_1 : z_2 \right\} $. Since $U_0 \cong \mathbb{C} ^2$, we have $\mathbb{C} \mathbb{P} ^2 \coloneqq \mathbb{C} ^2 \amalg \mathbb{C} ^1 \amalg \mathbb{C} ^0$. We can induct to obtain
\[
	\mathbb{C} \mathbb{P} ^n \coloneqq \coprod_{i=0}^{n} \mathbb{C} ^i
.\]
Hence $\mathbb{C} \mathbb{P} ^n$ admits this cell decomposition and is a CW complex with cells of even real dimension. Here we of course have not made the gluing relations explicit. Since cohomology of a CW complex comes by each level of the cells, and every other degree of cells is zero, the cohomology $H^{\bullet} (\mathbb{C} \mathbb{P} ^n;\mathbb{C} )\cong \mathbb{C} / 0 \cong  \mathbb{C} $. Oh and also $\mathbb{C} \mathbb{P} ^n \cong S^{n+1} /S^1$ a diffeomorphism.

Submanifolds (more generally, subvarieties) of $\mathbb{C} \mathbb{P} ^n$ admit descriptions as \emph{homogeneous polynomials}. A polynomial $f$ in $n+1$ variables is \emph{homogeneous of degree $d$} if $f(\lambda x)=\lambda ^df(x)$, $x\in\mathbb{C} ^{n+1} $. It follows that $f$ is a sum of monomials of degree $d$.
